\section{Exact Split-Zero carriers and the local spectral bridge}
\label{sec:split-zero-carriers}

Throughout this section, semiring homomorphisms preserve additive zero and multiplicative unit. Let $R$ be a nonzero commutative ring and $\mathbb B=(\{0,1\},\vee,\wedge)$. Define
\[
G(R)=R\sqcup\{\tau\},\qquad e=0_R\in R\subset G(R).
\]
Addition and multiplication of supported elements are those of $R$, including when their result is zero. The additional element satisfies
\[
\tau+r=r+\tau=r,\quad \tau+\tau=\tau,\quad \tau r=r\tau=\tau,\quad \tau^2=\tau .
\]
Thus the semiring zero is $\tau$, and $e\ne\tau$.

\subsection{Scalar support and independent channel amplitudes}

\begin{proposition}[Pair carrier and scalar reflection]
\label{prop:sz-pair-reflection}
There is a semiring isomorphism
\[
\eta_R:G(R)\longrightarrow \{(0,0)\}\cup(R\times\{1\})\subset R\times\mathbb B,
\qquad \tau\mapsto(0,0),\quad r\mapsto(r,1).
\]
The coordinate maps
\[
p_R:G(R)\to R,\quad p_R(\tau)=0,\quad p_R(r)=r,
\]
\[
\chi_R:G(R)\to\mathbb B,\quad \chi_R(\tau)=0,\quad\chi_R(r)=1
\]
are semiring homomorphisms. The map $p_R$ is universal among semiring homomorphisms from $G(R)$ to a ring.
\end{proposition}
\begin{proof}
In the product semiring, two supported pairs have sum $(r+s,1)$ and product $(rs,1)$. These identities retain support when the amplitude vanishes. The pair $(0,0)$ is an additive identity and an absorbing multiplicative zero. This proves closure and all operation identities, and the displayed map is bijective. Its coordinate projections give $p_R,\chi_R$. If $f:G(R)\to A$ has ring target, then $e+e=e$ gives $f(e)+f(e)=f(e)$, hence $f(e)=0_A$. Its restriction to the supported copy of $R$ is therefore a unital ring homomorphism $h:R\to A$. The equality $f=h p_R$ holds on $R$ and on $\tau$, and determines $h$. Conversely every such $h$ gives $h p_R$.
\end{proof}

The larger carrier $R\times\mathbb B$ contains $(1,0)$; its exact relation to $G(R)$ is the inclusion above. The pair presentation requires $R\times\{1\}$.

For a nontrivial bounded distributive lattice $L$, define
\[
G_L(R)=\{(0,\lambda):\lambda\in L\}\cup\{(r,1_L):r\in R\}\subset R\times L,
\]
\[
(r,\lambda)+(s,\mu)=(r+s,\lambda\vee\mu),\qquad
(r,\lambda)(s,\mu)=(rs,\lambda\wedge\mu).
\]
Write $z_\lambda=(0,\lambda)$, $\tau_L=z_{0_L}$, and $e_L=z_{1_L}$.

\begin{proposition}[Lattice reflection with its original coefficient ring]
\label{prop:sz-lattice-reflection}
The displayed carrier is a commutative semiring. The projections
\[
p_L:G_L(R)\to R,\qquad \chi_L:G_L(R)\to L
\]
are respectively the universal ring reflection and the support homomorphism. In particular the amplitude codomain of $G_{\mathbb B^2}(\mathbb C)$ is $\mathbb C$.
\end{proposition}
\begin{proof}
A summand with nonzero amplitude has top support, so its sum has top support. A factor with lower support has zero amplitude, so a product of lower support has zero amplitude. These facts prove closure. All laws restrict from the product semiring; distributivity uses that of $L$. Both projections preserve operations, zero and unit. If $f:G_L(R)\to A$ has ring target, $e_L+z_\lambda=e_L$ gives $f(z_\lambda)=0$ by cancellation in $A$. Restriction to the supported copy of $R$ is a ring homomorphism $h$ and $f=h p_L$. Conversely each $h$ gives this factorization. Surjectivity of $p_L$ gives uniqueness.
\end{proof}

Let $j:\mathbb C\to G(\mathbb C)$ include the supported copy. It preserves addition, multiplication and unit; it takes ordinary zero to $e$, whereas the semiring zero is $\tau$. Put
\[
M=G(\mathbb C)^2,\quad
p^{(2)}=(p_{\mathbb C},p_{\mathbb C}):M\to\mathbb C^2,\quad
\chi^{(2)}=(\chi_{\mathbb C},\chi_{\mathbb C}):M\to\mathbb B^2.
\]
For $\epsilon\in\mathbb B$ let $u_0=\tau$, $u_1=e$.

\begin{proposition}[Exact scalar-lattice and two-amplitude embeddings]
\label{prop:sz-two-channel-embeddings}
There are injective semiring homomorphisms
\[
b:G_{\mathbb B^2}(\mathbb C)\hookrightarrow M,\qquad
B:G_{\mathbb B^2}(\mathbb C^2)\hookrightarrow M
\]
defined by
\[
\begin{aligned}
b(c,(1,1))&=(j(c),j(c)),\\
b(0,(\lambda_1,\lambda_2))&=(u_{\lambda_1},u_{\lambda_2}),\\
B((a,b),(1,1))&=(j(a),j(b)),\\
B((0,0),(\lambda_1,\lambda_2))&=(u_{\lambda_1},u_{\lambda_2}).
\end{aligned}
\]
Their amplitude and support identities are
\[
p^{(2)}b=\Delta_{\mathbb C}p_{\mathbb B^2},\quad
\chi^{(2)}b=\chi_{\mathbb B^2},\qquad
p^{(2)}B=p_{\mathbb B^2},\quad
\chi^{(2)}B=\chi_{\mathbb B^2},
\]
where $\Delta_{\mathbb C}(c)=(c,c)$.
\end{proposition}
\begin{proof}
The pair presentation identifies $M$ with
\[
\left\{((a,b),(\lambda_1,\lambda_2))\in\mathbb C^2\times\mathbb B^2:
\lambda_1=0\Rightarrow a=0,\quad\lambda_2=0\Rightarrow b=0\right\},
\]
with componentwise operations. The map $b$ restricts the product semiring map $(c,\lambda)\mapsto((c,c),\lambda)$. Lower support in its source forces $c=0$, so it lands in the displayed subset. It retains the amplitude and both support bits and is injective. The map $B$ restricts the identity coordinate map on $\mathbb C^2\times\mathbb B^2$; lower support in its source forces both amplitudes to vanish, so it too lands in $M$ and is injective. The prescriptions agree at the top-support zero. Both maps preserve zero and unit. Projection proves all four identities.
\end{proof}

The element $(j(1),\tau)\in M$ is outside both images. There is no function $r:M\to G_{\mathbb B^2}(\mathbb C^2)$ satisfying both $p_{\mathbb B^2}r=p^{(2)}$ and $\chi_{\mathbb B^2}r=\chi^{(2)}$: it would send that element to $((1,0),(1,0))$, outside its target carrier. This identifies the precise obstruction to a reverse map preserving both projections.

The synchronized carrier has the exact map
\[
G(\mathbb C^2)\cong G(\mathbb C)\times_{\mathbb B}G(\mathbb C)\hookrightarrow M.
\]
It takes $\tau$ to $(\tau,\tau)$ and a supported $(a,b)$ to $(j(a),j(b))$. Equality of the two support characters leaves exactly those cases, proving bijectivity onto the fibre product. Operations agree componentwise. Its image has masks $(0,0),(1,1)$, while the other masks in $M$ carry independent mixed amplitudes.

The map $p^{(2)}$ is the universal ring reflection of $M$. Indeed, $(e,e)$ is additively idempotent, so every ring-target map kills it. Adding $(e,e)$ to $x\in M$ gives the fully supported element of amplitude $p^{(2)}x$. The map is consequently determined by a ring homomorphism from the supported copy of $\mathbb C^2$, proving the universal property. A semiring isomorphism $G_{\mathbb B^2}(\mathbb C)\cong M$ would induce a ring isomorphism $\mathbb C\cong\mathbb C^2$. This is impossible because $(1,0)$ is a nontrivial idempotent of $\mathbb C^2$, and $\mathbb C$ has only $0,1$. The positive relation between the original carriers is $b$, with diagonal amplitude projection.

\subsection{Tesserine sectors and a supported critical fixed locus}

Reserve $\tau$ for split absence and use $\kappa$ for conjugation. Let
\[
\sigma(s)=1-s,\quad\kappa(s)=\bar s,\quad K_4=\{1,\sigma,\kappa,\sigma\kappa\}.
\]
The real tesserine algebra and its complexification are
\[
A_{\mathbb R}=\mathbb R[K_4]
=\mathbb R[\sigma,\kappa]/(\sigma^2-1,\kappa^2-1),\qquad
A_{\mathbb C}=\mathbb C\otimes_{\mathbb R}A_{\mathbb R}.
\]
The polynomial algebra in this presentation is commutative.

\begin{proposition}[Character algebra and quartet coordinates]
\label{prop:sz-quartet-coordinates}
For $\epsilon,\delta\in\{1,-1\}$ let
\[
E_{\epsilon\delta}=\tfrac14(1+\epsilon\sigma)(1+\delta\kappa).
\]
Character evaluation gives an algebra isomorphism
\[
\mathcal F:A_{\mathbb C}\to\mathbb C^4,\qquad
x\mapsto(\operatorname{ev}_{\epsilon,\delta}(x))_{\epsilon,\delta},
\]
with inverse $(d_{\epsilon\delta})\mapsto\sum d_{\epsilon\delta}E_{\epsilon\delta}$.
For the exact quartet element
\[
Q(s)=s\cdot1+(1-s)\sigma+\bar s\,\kappa+(1-\bar s)\sigma\kappa,\qquad s=\beta+i\gamma,
\]
the four character values are
\[
c_{++}=2,\qquad c_{+-}=0,\qquad c_{-+}=2(2\beta-1),\qquad c_{--}=4i\gamma.
\]
Projection to the last two specified characters is a unital algebra map $\pi:A_{\mathbb C}\to\mathbb C^2$. The composite $\pi Q$ is a real affine bijection
\[
\mathbb C\to\mathbb R\times i\mathbb R,\qquad
s\mapsto(a,b)=(c_{-+}(s),c_{--}(s)),\qquad s=\tfrac12+\tfrac14(a+b).
\]
\end{proposition}
\begin{proof}
Evaluation of $E_{\epsilon\delta}$ at $(\epsilon',\delta')$ is $\tfrac14(1+\epsilon\epsilon')(1+\delta\delta')$. It is one when both signs agree and zero otherwise. This proves the two inverse formulas and all orthogonal idempotent identities, retaining the factors of four. Character evaluation preserves the algebra operations. Directly,
\[
\begin{aligned}
c_{++}&=s+(1-s)+\bar s+(1-\bar s)=2,\\
c_{+-}&=s+(1-s)-\bar s-(1-\bar s)=0,\\
c_{-+}&=s-(1-s)+\bar s-(1-\bar s)=2(s+\bar s)-2=4\beta-2,\\
c_{--}&=s-(1-s)-\bar s+(1-\bar s)=2(s-\bar s)=4i\gamma.
\end{aligned}
\]
The proposed inverse recovers $\beta+i\gamma$. Conversely every real $a$ and purely imaginary $b$ is recovered by substituting that inverse into the two character expressions.
\end{proof}

The quartet image has real affine dimension two. Its real linear span has dimension three, spanned by $E_{++},E_{-+},iE_{--}$, retaining the constant character coordinate. The real tesserine algebra has dimension four. The maps $Q,\mathcal F,\pi$ and the displayed affine inverse specify the exact relationship of these objects.

\begin{proposition}[Supported sector lift and fixed locus]
\label{prop:sz-critical-fixed-locus}
Define
\[
\Psi:\mathbb C\to M,\qquad\Psi(s)=(j(c_{-+}(s)),j(c_{--}(s))).
\]
It has $\chi^{(2)}\Psi(s)=(1,1)$ for all $s$. The $G(\mathbb C)$-linear automorphisms
\[
T_\sigma(x,y)=((-1)x,(-1)y),\qquad T_\kappa(x,y)=(x,(-1)y)
\]
of the free semimodule $M$ give
\[
\Psi(\sigma s)=T_\sigma\Psi(s),\qquad\Psi(\kappa s)=T_\kappa\Psi(s).
\]
Here $-1=j(-1)$ is a supported unit. Their product satisfies
\[
\operatorname{Fix}(T_\sigma T_\kappa)=\{\tau,e\}\times G(\mathbb C),\qquad
\Psi^{-1}\operatorname{Fix}(T_\sigma T_\kappa)=\{s:\Re s=\tfrac12\}.
\]
\end{proposition}
\begin{proof}
The sector formulas give $(a,b)(1-s)=(-a,-b)$ and $(a,b)(\bar s)=(a,-b)$. Multiplication by $j(-1)$ negates supported amplitudes and fixes $\tau$. It preserves addition, commutes with scalar multiplication and squares to identity. This proves the automorphism and equivariance statements. Its fixed points are $\tau$ and supported $j(c)$ with $c=-c$. Since $2\ne0$ in $\mathbb C$, the latter is precisely $e$. The product action is $((-1)x,y)$, proving its fixed locus. On the image of $\Psi$ the first coordinate is supported and is fixed exactly when $4\beta-2=0$.
\end{proof}

The scalar-lattice channel atoms have images $b(z_{10})=(e,\tau)$ and $b(z_{01})=(\tau,e)$. The separate support swap has fixed masks $(0,0),(1,1)$. Since $\chi^{(2)}\Psi$ is constantly $(1,1)$, it cannot detect the critical line by that swap. The support action induced by $T_\sigma T_\kappa$ is identity; its amplitude action is $(-a,b)$. The preceding proposition supplies the exact fixed-locus test with the first-channel supported zero retained.

\subsection{Exact sector, holonomy and positive-moment maps}

Fix a prime $p$, $\ell_p=\log p>0$, and $u\in S^1$. For an arithmetic packet, take $u=\psi(\operatorname{Frob}_p)$ with $L/\mathbb Q$ finite abelian, $p$ unramified, and $\psi:\operatorname{Gal}(L/\mathbb Q)\to S^1$ a character. All powers use $p^w=\exp(w\ell_p)$ with the positive real logarithm.

\begin{proposition}[Sector-to-holonomy map with constants retained]
\label{prop:sz-sector-holonomy}
The analytic function
\[
H_{p,u}:\mathbb C^2\to\mathbb C^\times,\qquad
H_{p,u}(a,b)=u\exp\!\left(\frac{\ell_p}{4}(a+b)\right)
\]
gives a function $H_{p,u}p^{(2)}$ on $M$. On the quartet image,
\[
H_{p,u}(c_{-+}(s),c_{--}(s))=u p^{s-1/2}=:z_{p,u}(s),
\]
\[
\frac{\log|z_{p,u}(s)|}{\ell_p}
=\frac{c_{-+}(s)}4=\Re s-\tfrac12,\qquad
|z_{p,u}(s)|^2-1=\exp\!\left(\frac{\ell_p}{2}c_{-+}(s)\right)-1.
\]
\end{proposition}
\begin{proof}
The exact inverse for the sector coordinates gives $s-\tfrac12=(a+b)/4$. Substitution proves the exponential identity. On that image $a$ is real and $b$ is purely imaginary. Since $|u|=1$, taking absolute values yields $\exp(\ell_p a/4)$. Taking its real logarithm and, separately, its square proves the two remaining identities.
\end{proof}

The radial exponent is $c_{-+}/4$; the exponent $c_{-+}/2$ belongs to the squared modulus. The map $H_{p,u}$ is used with its analytic-function type, without any assertion that it preserves semiring addition.

\begin{proposition}[Rank-one packet]
\label{prop:sz-rank-one-packet}
For $z\in\mathbb C$ define
\[
v_2(z)=\binom1z,\quad T_2(z)=v_2(z)v_2(z)^*
=\begin{pmatrix}1&\bar z\\z&|z|^2\end{pmatrix},\quad P_z(w)=w-\bar z.
\]
Then $T_2(z)$ is Hermitian positive semidefinite of rank one and
\[
\ker T_2(z)=\mathbb C(-\bar z,1)^\mathsf T.
\]
It is Toeplitz exactly when $|z|=1$, equivalently when the unique root of $P_z$ lies on $S^1$. For $z=z_{p,u}(s)$ this is equivalent to $\Re s=\tfrac12$.
\end{proposition}
\begin{proof}
For $x\in\mathbb C^2$, $x^*T_2(z)x=|v_2(z)^*x|^2\ge0$. The vector $v_2(z)$ is nonzero and hence its outer product has rank one. The adjoint kernel equation is $x_0+\bar z x_1=0$, giving the displayed space. A two-by-two matrix is Toeplitz exactly when its diagonal entries agree, here $1=|z|^2$. The root $\bar z$ has modulus $|z|$. The final equivalence follows from the preceding proposition and $\ell_p>0$.
\end{proof}

\begin{theorem}[Positive finite packets and their exact diagonal defect]
\label{thm:sz-finite-packet}
Let $r\ge1$, $\alpha_j>0$, and $z_j\in\mathbb C$. For $n\ge1$ put
\[
v_n(z)=(1,z,\ldots,z^{n-1})^\mathsf T,\qquad
M_n=\sum_{j=1}^r\alpha_jv_n(z_j)v_n(z_j)^* .
\]
The matrix $M_n$ is Hermitian positive semidefinite. For $n\ge3$, it is Toeplitz exactly when every $|z_j|=1$. With matrix indices starting at zero,
\[
\mathcal D(M_n):=(M_n)_{22}-2(M_n)_{11}+(M_n)_{00}
=\sum_{j=1}^r\alpha_j(|z_j|^2-1)^2.
\]
For these packets and $n\ge3$, vanishing of $\mathcal D$ already characterizes Toeplitzness.
If the $z_j$ are pairwise distinct and $r\le n$, the rank is $r$. If $r<n$, coefficient vectors identify the kernel with
\[
\left(\prod_{j=1}^r(w-\bar z_j)\right)\mathbb C[w]_{\le n-1-r}
\subset\mathbb C[w]_{\le n-1}.
\]
In particular $r=n-1$ gives corank one and exactly the conjugated roots $\bar z_1,\ldots,\bar z_{n-1}$.
\end{theorem}
\begin{proof}
The identity
\[
x^*M_nx=\sum_j\alpha_j|v_n(z_j)^*x|^2
\]
proves positivity and shows that its kernel is the simultaneous kernel of all $v_n(z_j)^*$. The first three diagonal entries are $\sum\alpha_j$, $\sum\alpha_j|z_j|^2$, and $\sum\alpha_j|z_j|^4$. Their displayed linear combination gives the defect without rescaling any weight. Toeplitzness equates these entries, forcing zero defect. Every weight is positive and every square nonnegative, hence all $|z_j|=1$. Conversely that condition gives $(M_n)_{k\ell}=\sum_j\alpha_jz_j^{k-\ell}$, a Toeplitz matrix. The same argument proves the zero-defect characterization.

Let $V$ have columns $v_n(z_j)$. Its first $r$ rows have Vandermonde determinant $\prod_{i<j}(z_j-z_i)\ne0$, so $V$ has rank $r$. The established kernel identity is $\ker M_n=\ker V^*$, proving the rank. For $P_x(w)=\sum_{k=0}^{n-1}x_kw^k$, the kernel equations say $P_x(\bar z_j)=0$. Distinct linear factors divide $P_x$ exactly when these equations hold: division by one factor preserves the other root equations since that factor is nonzero at those roots. This proves the product and degree bound. At $r=n-1$, the remaining factor is a scalar; for a nonzero kernel vector it is nonzero, so there are no further roots.
\end{proof}

The size bound is sharp for general positive packets. At $n=2$, the data $(\alpha_1,\alpha_2)=(4,1)$ and $(z_1,z_2)=(1/2,2)$ give
\[
M_2=4\begin{pmatrix}1&1/2\\1/2&1/4\end{pmatrix}
+\begin{pmatrix}1&2\\2&4\end{pmatrix}
=\begin{pmatrix}5&4\\4&5\end{pmatrix}.
\]
This is positive definite and Toeplitz, while neither atom is on $S^1$. At size three the diagonal is $(5,5,65/4)$ and the defect is $45/4$. One atom retains the two-by-two criterion. Zero-weight atoms are invisible; repeated atoms combine their weights before applying the distinct-column rank count.

The exact sector defect is therefore
\[
\mathcal D(M_n)=\sum_{j=1}^r\alpha_j
\left[\exp\!\left(\frac{\log p_j}{2}c_{-+}(s_j)\right)-1\right]^2 .
\]
The function chain
\[
M\xrightarrow{p^{(2)}}\mathbb C^2
\xrightarrow{H_{p,u}}\mathbb C^\times
\xrightarrow{z\mapsto v_n(z)v_n(z)^*}\operatorname{Herm}_{\ge0}(n)
\]
has explicit maps at every step. Positive finite summation supplies the preceding theorem.

\subsection{The theta-packet form and its scalar comparison}

\begin{proposition}[Eigenpair form quotient and sector channel]
\label{prop:sz-theta-form-quotient}
For every complex $d$-by-$d$ matrix $A$, Hermitian matrix $G$, and nonzero eigenvector $v$ with $Av=\rho v$, let $W=A^*G+GA-G$. Then
\[
v^*Wv=(2\Re\rho-1)v^*Gv.
\]
On the domain of these data where $v^*Gv\ne0$, the quotient
\[
\omega(A,G,v,\rho)=\frac{v^*Wv}{v^*Gv}
=2\Re\rho-1=\frac{c_{-+}(\rho)}2
\]
is real. Its associated holonomy satisfies
\[
|z_{p,u}(\rho)|^2=\exp(\ell_p\omega),\qquad
z_{p,u}(\rho)=u\exp\!\left(\frac{\ell_p}{4}
(c_{-+}(\rho)+c_{--}(\rho))\right).
\]
\end{proposition}
\begin{proof}
Since $(Av)^*=\bar\rho\,v^*$,
\[
v^*(A^*G+GA-G)v
=(Av)^*Gv+v^*G(Av)-v^*Gv
=(\bar\rho+\rho-1)v^*Gv .
\]
This proves the unquotiented identity even when the right factor is zero. Division on the stated domain gives $\omega$, and $c_{-+}(\rho)=4\Re\rho-2$ gives the factor $1/2$. The holonomy identities follow from the exact sector exponential above.
\end{proof}

This is the typed map from an eigenpair of the form packet to its Split-Zero sector and local moment matrix. The moment matrix is positive for every parameter, so that positivity cannot supply a further zero-defect conclusion. For $\rho=1$, $u=1$,
\[
z=\sqrt p,\qquad T_2(z)=\begin{pmatrix}1&\sqrt p\\\sqrt p&p\end{pmatrix}\ge0,
\]
and its diagonal entries differ. On the form side $A=(1)$, $G=(1)$, $v=(1)$ give $W=(1)$ and $\omega=1$; the exact maps take these data to the same off-critical packet. Any arithmetic identity forcing the actual family of $W$ to vanish on its eigenvectors must therefore be established from that family's construction. These maps retain that requirement rather than replacing it by Gram positivity.

\subsection{Projective zeta at all zeros and at its pole}

A pair in $G(\mathbb C)^2$ is unimodular exactly when at least one amplitude is nonzero. Such a coordinate is a supported unit. If both amplitudes are zero, every linear combination has zero amplitude and cannot be the unit. Let $\mathbb P_{\rm um}^1(G(\mathbb C))$ denote unimodular pairs modulo common unit multiplication.

\begin{proposition}[Projective amplitude section and reciprocal]
\label{prop:sz-projective-section}
The map
\[
p_*:\mathbb P_{\rm um}^1(G(\mathbb C))\to\mathbb P^1(\mathbb C),
\quad[x:y]\mapsto[p(x):p(y)]
\]
has an injective section
\[
\iota:\mathbb P^1(\mathbb C)\to\mathbb P_{\rm um}^1(G(\mathbb C)),
\quad[a:b]\mapsto[j(a):j(b)],
\]
whose image is the full-support component. With
\[
\mathcal R:G(\mathbb C)\to\mathbb P_{\rm um}^1(G(\mathbb C)),
\quad x\mapsto[1:x],\qquad
r:\mathbb C\to\mathbb P^1(\mathbb C),\quad c\mapsto[1:c],
\]
one has $p_*\mathcal R=r p_{\mathbb C}$ and $\mathcal Rj=\iota r$.
\end{proposition}
\begin{proof}
Unimodularity ensures the projected pair is nonzero. A common supported unit projects to a nonzero complex scalar, so $p_*$ is well defined. A nonzero ordinary pair becomes a unimodular supported pair; changing its representative lifts to common supported unit multiplication. Thus $\iota$ is well defined, and $p_*\iota$ is identity. Every full-support pair is the supported lift of its amplitude pair. The pair $(1,x)$ is always unimodular, and projection of its coordinates proves both remaining identities.
\end{proof}

In particular $\mathcal R(e)=[1:e]=\infty_e$ and $\mathcal R(\tau)=[1:\tau]=\infty_\tau$. Both project to $[1:0]$, while their Boolean projective supports are $[1:1]$ and $[1:0]$. Similarly $[e:1]$ and $[\tau:1]$ project to the ordinary zero while retaining different supports.

\begin{proposition}[Two exact charts for the supported zeta map]
\label{prop:sz-zeta-charts}
For the meromorphic zeta function with its simple pole of residue one at $1$, the full-support lift is
\[
Z_{\rm sp}:\mathbb C\to\mathbb P_{\rm um}^1(G(\mathbb C)),\qquad
Z_{\rm sp}(s)=
\begin{cases}
[j(\zeta(s)):1],&s\ne1,\\
[1:e],&s=1 .
\end{cases}
\]
It equals $\iota\widehat\zeta$. At every zeta zero $\rho$ its value is $[e:1]$. Near the pole it has the chart
\[
Z_{\rm sp}(s)=[1:j(h(s))],\qquad h(s)=1/\zeta(s),\quad h(1)=0,\quad h'(1)=1.
\]
The reciprocal is a finite-amplitude coordinate away from the zeta zeros, including its extension at $1$. At a zero of order $m$, it has a pole of order $m$.
\end{proposition}
\begin{proof}
Put $t=s-1$ and write $\zeta(1+t)=t^{-1}+g(t)$ with $g$ holomorphic near zero. Then $h(1+t)=t/(1+t g(t))$. Its denominator has value one at zero, so this is holomorphic there with value zero and derivative one. On the overlap where $\zeta\ne0$, common multiplication by $j(\zeta(s))$ gives $[1:j(h(s))]=[j(\zeta(s)):1]$. This proves the transition and pole value. The first chart directly gives $[e:1]$ at each zero. At a zero of order $m$, write $\zeta(\rho+t)=t^m u(t)$ with $u(0)\ne0$. Its reciprocal is $t^{-m}u(t)^{-1}$, proving its pole order; the first chart retains the original order $m$.
\end{proof}

The swap $J_*[x:y]=[y:x]$ is induced by an invertible permutation matrix and commutes with $\iota,p_*$. Hence $J_*Z_{\rm sp}$ is the global reciprocal projective map, sending zeta zeros to $\infty_e$ and the pole to $[e:1]$. Analytic statements use the complex structure transported to the full-support component by $\iota$; no topology on the extra support components is needed.

\subsection{Matrix packets and their full relation module}

\begin{theorem}[Block packet with explicit hypotheses]
\label{thm:sz-block-packet}
Let $V$ be a complex Hilbert space of finite dimension $d>0$ with fixed positive-definite Hermitian inner product. Let $g\in\operatorname{GL}(V)$ and $m\ge1$. Use that metric and its orthogonal direct sums for all adjoints. Define
\[
W_m(g):V\to V^{m+1},\quad v\mapsto(v,gv,\ldots,g^mv),\qquad
\mathbb T_m(g)=W_m(g)W_m(g)^*,
\]
\[
D_m(g):V^m\to V^{m+1},\qquad
\begin{cases}
(D_m(g)y)_0=-g^*y_1,\\
(D_m(g)y)_k=y_k-g^*y_{k+1}&(1\le k\le m-1),\\
(D_m(g)y)_m=y_m .
\end{cases}
\]
The matrix $\mathbb T_m(g)$ is positive semidefinite and
\[
\ker\mathbb T_m(g)=\operatorname{im}D_m(g),\qquad
\operatorname{rank}\mathbb T_m(g)=d,\qquad
\operatorname{corank}\mathbb T_m(g)=md .
\]
It is block Toeplitz exactly when $g$ is unitary for the fixed metric.
\end{theorem}
\begin{proof}
The Gram factorization gives positivity and $\ker\mathbb T_m=\ker W_m^*$. The first component of $W_mv$ is $v$, so $W_m$ is injective and the kernel dimension is $(m+1)d-d=md$. The adjoint is
\[
W_m^*(v_0,\ldots,v_m)=\sum_{k=0}^m(g^*)^k v_k.
\]
Substitution of $D_m(y)$ gives
\[
-g^*y_1+\sum_{k=1}^{m-1}
\big((g^*)^ky_k-(g^*)^{k+1}y_{k+1}\big)+(g^*)^my_m=0.
\]
Thus its image lies in the kernel. If $D_m(y)=0$, the last component gives $y_m=0$, and backward substitution gives each preceding $y_k=0$. Its image dimension is $md$, proving equality with the kernel. The block of indices $j,k$ is $g^j(g^*)^k$. For unitary $g$ this is $g^{j-k}$, hence block Toeplitz. Conversely block Toeplitzness equates blocks $(0,0),(1,1)$, giving $I=gg^*$. Invertibility yields $g^*=g^{-1}$ and unitarity.
\end{proof}

At $m=0$ the only block is $I$, so the bound $m\ge1$ is essential. At $m=1$, the full kernel is $\{(-g^*y,y):y\in V\}$; its vector polynomials are $(wI-g^*)y$. Taking determinant gives the exact scalar map
\[
wI-g^*\longmapsto\det(wI-g^*)=\overline{\det(\bar wI-g)}.
\]
The equality follows by conjugating the determinant expansion and transposing. It retains conjugate eigenvalues and multiplicities. For the standard metric and $J_t=\left(\begin{smallmatrix}1&t\\0&1\end{smallmatrix}\right)$ with $t>0$, the polynomial is $(w-1)^2$ but
\[
J_tJ_t^*=\begin{pmatrix}1+t^2&t\\t&1\end{pmatrix}\ne I.
\]
Its block packet consequently fails Toeplitzness, although the determinant roots lie on the unit circle. The kernel remains $\{(-J_t^*y,y):y\in\mathbb C^2\}$, retaining the matrix information that determinant forgets.
